\chapter{Semantic Analyzer}
\label{ch:semantics}

\section{Overview}

The semantic pass walks the AST produced by the parser and performs the following operations:
\begin{itemize}
    \item{\textbf{Name resolution}}: identifiers are resolved to their declaration sites.
    \item{\textbf{Type checking}}: every expression is assigned a \texttt{ZType}; incompatible opverands produce errors.
    \item{\textbf{Scope management}}: a stack of symbol tables tracks visibility of variables, functions and types.
    \item{\textbf{Generic instantiation}}: generic strucs and fuunctions are monomorphized on demand (not yet implemented).
    \item{\textbf{Capability checking}}: \zn{with T x} parameters are validated in their enclosing scope.
    \item{\textbf{Unused symbol warnings}}: variables, functions and structs that are never referenced emit warnings (suppressed by CLI flags).
\end{itemize}

The pass is conceptually split in two phases. The first phase (\texttt{discoverGlobalScope}) registers every top-level declaration (funcs, structs, enums, typedefs, namespaces, globals, foreign funcs) into the global scope, so that modules can refer to each other without ordering constraints. The second phase (\texttt{analyze}) walks each module and type-checks function bodies, statements and expressions.

\section{Symbol Tables}
A \texttt{ZSymTable} owns the global scope and the per-module scopes.
Each \texttt{ZScope} carries:
\begin{itemize}
    \item a \texttt{symbols} vector of \texttt{ZSymbol *} declared in the scope.
    \item a \texttt{seen} hashset used to detect duplicate names (\_ is ignored).
    \item a \texttt{capabilities} vector (see \ref{sec:capability_checking}).
    \item a pointer to its parent scope, used for lookup.
\end{itemize}

\texttt{resolve} walks the scope chain from the current scope upwards. On a successful match it bumps the symbol's \texttt{useCount} (the basis of the unused symbol warnings).
The placeholder identifier \zn{_} is rejected here so it cannot be referenced as a value.

\begin{lstlisting}[style=zinc]
sum :: i32 (i32 a, i32 b) {
    return a + b
}
main :: i32 () {
    i32 a = 10
    i32 b = 10
    res := sum(a, b)
    return 0
}
\end{lstlisting}

In the example above, \texttt{sum} and \texttt{main} are registered in the file (module) scope during phase one. When \texttt{main}'s body is analyzed in phase two, \texttt{beginScope} opens a fresh child scope for its parameters and locals; \texttt{endScope} pops it on exit and triggers the unused symbol check for that frame.

\section{Type Checking}
Every expression node is passed through \texttt{resolveType}, which dispatches on \texttt{ZNodeType} and caches the result in \texttt{node->resolved} so each node is typed at most once.
I created a function \texttt{typesCompatible} which checks if two types can be compatible and returns the promoted type.
This function is used everywhere to permit implici cast. For example a functino accepts \zn{i32} bit the user passes \zn{i8} without that function the compiler will reject it because their types are not the same.

\subsection{Binary expressions}

\texttt{resolveBinary} resolves both operants then asks \texttt{typesCompatible} for the common type. The rules are:

\begin{itemize}
    \item Two pointers (or pointer vs \zn{none}) compare equal and yield the pointer type.
    \item Identical types return as-is.
    \item For two primitives the type with the higher rank wins (\zn{char} < \zn{i8}/\zn{u8} < \zn{i16}/\zn{u16} < \zn{i32}/\zn{u32} < \zn{i64}/\zn{u64} < \zn{f32} < \zn{f64}).
    \item A float on either side promotes the result to the wider float.
    \item Mixed signed/unsigned: if the signed operand has the higher rank, it wins; if the raks are equal, the result is promoted to the next signed rank. Promotion past \zn{i64} is refused and a warning sugggests an explicit cast.
    \item Bit-operators (\zn{&}, \zn{|}, \zn{^}, \zn{<<}, \zn{>>}) reject non-integer operands.
\end{itemize}

After promotion the operands are wrapped in \texttt{NODE\_CAST} nodes through \texttt{implicitCast} when needed.
Comparison and logical operators (\zn{==}, \zn{!=}, \zn{<}, \zn{>}, \zn{<=}, \zn{>=}, \zn{and}, \zn{or}) always yield \zn{u1}. Assignment (\zn{=}) requires the LHS to be an lvalue (identifier, subscript, member or pointer dereference) and yields the LHS type.

\subsection{Unary expressions}
\texttt{resolveUnary} handles threee operators specially:
\begin{itemize}
    \item \zn{&expr} wraps the operant type in \texttt{Z\_TYPE\_POINTER}.
    \item \zn{*expr} requires a pointer and yields the pointee.
    \item \zn{not expr} casts the operand to \zn{u1} and returns \zn{u1}.
\end{itemize}

\subsection{Pointer arithmetic}
Pointer arithmetic is not currently special-cased: pointer vs integer operations fall through the primitive compatibility check and are rejected. Pointer offsetting must be expressed through subscripts (\zn{p[i]}) which \texttt{resolveArrSubscript} accepts both arrays and pointer types (but arrays can't decay to pointer like normal C).

\section{Receiver Resolution}

Receiver functions (methods) are stored separately from free functions. \texttt{ZFuncTable} groups method declarations by their receiver's base type. This type is basically an hashmap where the key is the type and the value is a list of function declarations. Receiver methods live in \texttt{funcDef} while namespace-qualified static methods live in \texttt{staticFuncDef}.

\subsection{Method calls (\zn{expr.method(...)})}

\texttt{resolveMemberAccess} resolved the type of the LHS, dereferences pointer types (so \zn{*p}.method and \zn{p.method} behave the same), then:
\begin{enumerate}
    \item If thje base is a struct, it first searches the struct's fields.
    \item Failing that, it calls \texttt{resolveFuncCallEmbedded}, which walks the receiver-method tables for the type and, recursively, for any embedded fields. A method found via two distinct embedded paths is reported as a conflict.
    \item On success, \texttt{memberAccess.mangled} is set to the method's mangled name so codegen can emit a directo call and inject the self argument.
\end{enumerate}

\subsection{checkFunctionUsedAsValue} guards against passing a receiver method by name (e.g. \ as an argument): once \texttt{memberAccess.mangled} is set the node is boud to a method and cannot be treated asa value.

\subsection{Static Access (\zn{T::method})}

\texttt{resolveStaticAccess} looks up the base name and the property:
\begin{itemize}
    \item For a primitive or user type that has registered static methods, it returns the functino type of the matching entry.
    \item For an enum, it rewrites the node to \texttt{NODE\_ENUM\_LIT} and returns the enum type; codegen them emits a variant constructor.
    \item For a namespacem it searches the namespace block for a foreign declaration with the requested name and binds \texttt{staticAccess.mangled} to underlying C symbol.
\end{itemize}

\begin{note}
Function names are mangled, This means that the function name is renamed with specific rules:
For static functions the prefix is \_ZNM and for receiver functions the prefix is \_ZN.
Then the segments are attached with the length of the segment + the segment itself.
See the functions \texttt{mangler} and \texttt{manglerM} in \texttt{zmod.c}.
\end{note}

\section{Capability Checking}
Each \texttt{ZScope} carries a \texttt{capabilities} vector of \texttt{ZCapability} entries, indexed by type.
\texttt{pubCapability} either creates a fresh entry for an unseen type or appends to the existing one, so capabilities introduced by an enclosing \zn{capability} block remain visible inside nested scopes.

\texttt{resolveFuncCall} iterates the callee's declared capability list and runs \texttt{lookupScopedCapability}, which walks from the current scope outward loooking for a matching entry. The most recently pushed instance is selected and recorded in \texttt{call.capabilities} for codegen. A missing capability raises \texttt{Capability 'CapabilityName' not found in the current scope}.

\section{Struct, Enum and Embedding Checks}
In addition to per-expression typing, the pass enforces several structural invariants:
\begin{itemize}
    \item \texttt{analyzeStruct} resolved every field type and runs \texttt{isInfiniteSize} to reject by-value cycles such as \zn{struct Node \{ Node next \}}. The user is told to use a poibnter instead. Duplicate field names within a struct, and conflicts between a struct's own field and a field inherited through an embedded struct, are also reported.
    \item \texttt{analyzeEnum} resolves each variant's payload, checks for duplicate variant names and runs the same infinite-size check on payload fields.
    \item \texttt{analyzeBlock} flags statements after \zn{return}, \zn{break} or \zn{continue} as unreachable and truncates the block at that point.
\end{itemize}

